\documentclass[11pt,a4paper]{article}
\usepackage{fontspec}
\setmainfont{DejaVu Sans}[Scale=0.90]
\usepackage{amsmath,amssymb}
\usepackage[margin=2.4cm]{geometry}
\usepackage{booktabs,array,longtable,tabularx}
\usepackage{enumitem}
\usepackage{titlesec}
\usepackage{parskip}
\titlelabel{\thetitle.\quad}
\titleformat{\section}{\normalfont\large\bfseries}{\thesection.}{0.6em}{}
\titleformat{\subsection}{\normalfont\normalsize\bfseries}{\thesubsection}{0.6em}{}
\titlespacing*{\section}{0pt}{16pt}{6pt}
\titlespacing*{\subsection}{0pt}{12pt}{4pt}
\setlist[itemize]{leftmargin=1.4em,itemsep=1pt,topsep=3pt}
\newcolumntype{L}[1]{>{\raggedright\arraybackslash}p{#1}}
\begin{document}
\begin{center}
{\Large\bfseries Fiche de lecture --- Krapez \& Rigollet (2017)}\\[4pt]
{\small Commentaire sur Guo, Mandelis, Tolev \& Tang, \emph{J. Appl. Phys.} \textbf{121}, 095101 (2017)}\\[2pt]
{\small Version 1 --- 12 septembre 2026}
\end{center}

\vspace{4pt}

\section{Identification}

\begin{center}
\begin{tabularx}{\textwidth}{@{}lX@{}}
\toprule
Auteurs & J.-C. Krapez (ONERA, DOTA, Salon-de-Provence) ; F. Rigollet (Aix-Marseille Univ., CNRS, IUSTI, Marseille) \\
\addlinespace
Titre & \emph{Comment on ``Photothermal radiometry parametric identifiability theory for reliable
and unique nondestructive coating thickness and thermophysical measurements''} \\
\addlinespace
Article commenté & X. Guo, A. Mandelis, J. Tolev, K. Tang, \emph{Journal of Applied Physics}, volume 121,
numéro 9, article 095101, 2017 \\
\addlinespace
Accès & Dépôt arXiv, identifiant 1708.07362v1, août 2017 \\
\addlinespace
Nature & Note critique courte, sans données expérimentales propres, à visée de réfutation \\
\bottomrule
\end{tabularx}
\end{center}

\section{Objet de la note}

L'article commenté prétend établir que le signal de radiométrie photothermique obtenu en transfert
unidimensionnel sur un échantillon revêtu opaque permet d'identifier séparément trois paramètres du
revêtement : l'épaisseur, la diffusivité et la conductivité thermiques. La note établit la proposition
inverse : ces trois paramètres sont structurellement corrélés dans la configuration considérée, le
critère d'identifiabilité utilisé est fautif, et l'épaisseur qui en est déduite n'est pas exploitable.

La thèse est de nature structurelle et non statistique. Elle ne porte ni sur le niveau de bruit, ni sur
la qualité de l'ajustement, mais sur le rang de l'application paramètres $\rightarrow$ observable.

\section{Rappels physiques}

\subsection{Grandeurs thermiques et relations de dépendance}

\[ a = \frac{\lambda}{\rho c} , \qquad b = \sqrt{\lambda\,\rho c} , \qquad
\lambda = b\sqrt{a} , \qquad \rho c = \frac{b}{\sqrt{a}} . \]

Deux grandeurs parmi $\lambda$, $\rho c$, $a$, $b$ suffisent à décrire complètement un milieu homogène ;
les deux autres s'en déduisent. Cette redondance n'est pas en cause dans la controverse : le point en
litige porte sur le triplet $(L,\lambda,a)$, qui comporte une grandeur géométrique supplémentaire.

\subsection{Correspondance de notations}

L'article commenté et la note emploient un jeu de symboles propre à la métrologie photothermique
anglo-saxonne. La correspondance avec les conventions du projet est la suivante.

\begin{center}
\begin{tabularx}{\textwidth}{@{}llX@{}}
\toprule
\textbf{Note} & \textbf{Projet} & \textbf{Grandeur} \\
\midrule
$L_{2}$ & $e$ & Épaisseur du revêtement \\
$\alpha_{2}$ & $a$ & Diffusivité thermique du revêtement \\
$\kappa_{2}$ & $\lambda$ & Conductivité thermique du revêtement \\
$\rho_{2}C_{2}$ & $\rho c$ & Capacité thermique volumique du revêtement \\
$P_{2}$ & $b$ & Effusivité thermique du revêtement \\
$Q_{2}$ & $\xi_{1}$ & Racine du temps de transit thermique, $Q_{2}=L_{2}/\sqrt{\alpha_{2}}$ \\
$b_{32}$ & --- & Rapport d'effusivité substrat sur revêtement, $b_{32}=P_{3}/P_{2}$ \\
$p$ & $p$ & Variable de Laplace ; en régime harmonique, $p=2i\pi f$ \\
\bottomrule
\end{tabularx}
\end{center}

Le paramètre $Q_{2}$ coïncide exactement avec la coordonnée transformée de Liouville évaluée à la face
arrière du revêtement, puisque $\xi(z)=\int_{0}^{z}du/\sqrt{a(u)}$ se réduit à $e/\sqrt{a}$ en milieu
homogène. Le jeu $\{Q_{2},P_{2}\}$ est donc le couple $\{\xi_{1},b\}$ des conventions du projet.

\subsection{Régime harmonique et onde thermique}

En milieu semi-infini homogène soumis à un flux modulé de densité $\wp$,
\[ \theta_{0} = \frac{\wp}{b\sqrt{p}} , \qquad
|\theta_{0}| = \frac{\wp}{b\sqrt{2\pi f}} , \qquad \arg(\theta_{0}) = -45^{\circ} , \qquad
\mu = \sqrt{\frac{2a}{\omega}} . \]

L'effusivité gouverne seule l'amplitude de la réponse de surface ; la diffusivité n'intervient que par
la profondeur sondée $\mu$. Cette dissymétrie est à l'origine de toute la difficulté : l'épaisseur et la
diffusivité ne se manifestent qu'au travers du même temps caractéristique.

\subsection{Quadripôle de la couche}

La note rappelle que, en l'absence de sources internes, la contribution de la couche est entièrement
portée par le quadripôle
\[ \mathbf{M} =
\begin{pmatrix}
\cosh\!\left(Q_{2}\sqrt{p}\right) & \dfrac{\sinh\!\left(Q_{2}\sqrt{p}\right)}{P_{2}\sqrt{p}} \\[10pt]
P_{2}\sqrt{p}\,\sinh\!\left(Q_{2}\sqrt{p}\right) & \cosh\!\left(Q_{2}\sqrt{p}\right)
\end{pmatrix} ,
\qquad AD-BC = 1 . \]

Cette expression est identique, terme à terme, à celle retenue dans les conventions du projet pour le
mur homogène, à la section unitaire près : avec $k=\sqrt{p/a}$ et $\xi_{1}=e/\sqrt{a}$, les identités
$ke=\sqrt{p}\,\xi_{1}$ et $\lambda k = b\sqrt{p}$ assurent la coïncidence.

Seuls deux nombres caractérisent donc le revêtement dans le modèle direct : $Q_{2}$ et $P_{2}$. Les
trois grandeurs physiques $L_{2}$, $\alpha_{2}$, $\kappa_{2}$ n'y entrent pas séparément.

\subsection{Réponse de surface du système revêtement--substrat}

Avec $Z=1/(P_{3}\sqrt{p})$ l'impédance du substrat semi-infini et un échange en face avant nul,
\[ \theta_{0} = \wp\,\frac{AZ+B}{CZ+D}
= \frac{\wp}{P_{2}\sqrt{p}}\cdot
\frac{\cosh\!\left(Q_{2}\sqrt{p}\right)+b_{32}\sinh\!\left(Q_{2}\sqrt{p}\right)}
{b_{32}\cosh\!\left(Q_{2}\sqrt{p}\right)+\sinh\!\left(Q_{2}\sqrt{p}\right)} , \]
soit, sous forme d'ondes multiplement réfléchies,
\[ \theta_{0} = \frac{\wp}{P_{2}\sqrt{p}}\cdot
\frac{1+\Gamma\,e^{-2Q_{2}\sqrt{p}}}{1-\Gamma\,e^{-2Q_{2}\sqrt{p}}} ,
\qquad \Gamma = \frac{1-b_{32}}{1+b_{32}} . \]

Trois lectures immédiates s'en déduisent.

\begin{itemize}
\item Le revêtement n'intervient que par $P_{2}$, en facteur d'échelle, et par $Q_{2}$, dans
l'exponentielle. Aucune combinaison de mesures ne peut donc séparer $L_{2}$, $\alpha_{2}$ et
$\kappa_{2}$.
\item Pour $b_{32}=1$, $\Gamma$ s'annule : l'interface devient indétectable et $Q_{2}$ cesse d'être
identifiable, quelle que soit la méthode et quel que soit le régime, pulsé ou modulé.
\item Lorsque l'amplitude est normalisée par une référence, le facteur $P_{2}$ disparaît et l'observable
ne dépend plus que de $Q_{2}$ et $b_{32}$.
\end{itemize}

\section{Rappels mathématiques}

\subsection{Vecteurs de sensibilité et matrice d'information}

Soit $G(f_{i})$ l'observable, amplitude ou phase, mesurée aux fréquences $f_{i}$, et $\beta$ le vecteur
des paramètres recherchés. La matrice de sensibilité $S$ a pour éléments $S_{ij}=\partial G(f_{i})/
\partial\beta_{j}$. Dans l'approche des moindres carrés, la matrice d'information $S^{T}S$ doit être
inversée, et la matrice de covariance des estimations est proportionnelle à $\left(S^{T}S\right)^{-1}$,
dont la diagonale fournit les facteurs d'amplification de la variance du bruit.

\subsection{Dépendance linéaire et rang déficient}

Si les colonnes de $S$ satisfont une relation linéaire à coefficients non tous nuls, valable pour toutes
les fréquences, alors $\operatorname{rang}(S) < \dim\beta$, donc $\det\left(S^{T}S\right)=0$ et
l'inversion est impossible. La non-identifiabilité est alors dite \emph{inconditionnelle} : elle ne
dépend ni du nombre de fréquences, ni de leur choix, ni du niveau de bruit.

Une matrice $S^{T}S$ singulière peut néanmoins apparaître numériquement inversible lorsque les
sensibilités sont obtenues par différences finies : les erreurs d'arrondi résiduelles écartent le
déterminant de zéro. L'inversion aboutit alors formellement, sans que la corrélation ait disparu.

\subsection{Non-identifiabilité et conditionnement}

Deux notions doivent être distinguées.

\begin{center}
\begin{tabularx}{\textwidth}{@{}lXX@{}}
\toprule
& \textbf{Non-identifiabilité structurelle} & \textbf{Mauvais conditionnement} \\
\midrule
Origine & Rang déficient de l'application directe & Valeurs propres faibles mais non nulles \\
\addlinespace
Dépendance au bruit & Aucune & Déterminante \\
\addlinespace
Dépendance à la bande & Aucune & Déterminante \\
\addlinespace
Effet & Solution non unique, incertitude infinie & Solution unique, incertitude grande \\
\addlinespace
Remède & Information supplémentaire de nature différente & Élargissement de bande, réduction du bruit,
régularisation \\
\bottomrule
\end{tabularx}
\end{center}

La note se situe intégralement dans la première colonne. C'est la raison pour laquelle aucune
optimisation du plan d'expérience ne peut y remédier.

\subsection{Transposition entre domaines temporel et fréquentiel}

La note énonce un principe de transposition : lorsqu'une combinaison de paramètres est non identifiable
dans le domaine de Fourier pour tout jeu de fréquences, elle l'est également dans le domaine temporel
pour tout jeu d'instants, et réciproquement. Ce principe découle du fait que la dépendance
paramétrique s'exerce à travers le même modèle direct, la transformée intégrale n'étant qu'un changement
de représentation. Une dégénérescence de rang indépendante de $p$ est donc invariante par inversion de
la transformée.

\subsection{Fonctions homogènes et identité d'Euler}

Soit $G$ une fonction des paramètres, invariante sous l'action d'un groupe à un paramètre
$\beta_{j}\mapsto\mu^{\,n_{j}}\beta_{j}$. La dérivation de $G$ par rapport à $\mu$ en $\mu=1$ fournit
\[ \sum_{j} n_{j}\,\beta_{j}\,\frac{\partial G}{\partial\beta_{j}} = 0 . \]
Cette identité est la forme d'Euler associée au groupe. Toute invariance d'échelle du modèle direct se
traduit donc mécaniquement par une relation linéaire entre vecteurs de sensibilité, et réciproquement.
Cette lecture, développée à la section 5.3, donne à l'argument de la note une origine géométrique.

\section{Argument central}

\subsection{Réduction à deux paramètres}

Le modèle direct de l'article commenté peut être réécrit de sorte que les propriétés du revêtement
n'interviennent que par
\[ Q_{2} = \frac{L_{2}}{\sqrt{\alpha_{2}}} , \qquad P_{2} = \frac{\kappa_{2}}{\sqrt{\alpha_{2}}} . \]
La note observe que les auteurs de l'article commenté mentionnent eux-mêmes cette réduction, mais n'en
tirent pas la conséquence. Le formalisme des quadripôles la rend immédiate.

Les deux combinaisons obtenues par produit et par quotient ne lèvent pas l'indétermination :
\[ Q_{2}P_{2} = \rho_{2}C_{2}L_{2} , \qquad \frac{Q_{2}}{P_{2}} = \frac{L_{2}}{\kappa_{2}} , \]
soit la capacité thermique surfacique et la résistance thermique du revêtement. L'épaisseur y demeure
liée à une propriété thermophysique. Son extraction exige une mesure indépendante de $\alpha_{2}$,
$\rho_{2}C_{2}$ ou $\kappa_{2}$ ; la connaissance de la seule effusivité est insuffisante.

\subsection{Dépendance linéaire des vecteurs de sensibilité}

Puisque $G=G(Q_{2},P_{2})$, la règle de dérivation composée donne, pour toute fréquence $f_{i}$ :
\[
\begin{aligned}
\frac{\partial G(f_{i})}{\partial\alpha_{2}}
&= -\frac{L_{2}}{2\alpha_{2}^{3/2}}\,\frac{\partial G(f_{i})}{\partial Q_{2}}
   -\frac{\kappa_{2}}{2\alpha_{2}^{3/2}}\,\frac{\partial G(f_{i})}{\partial P_{2}} \\[2pt]
\frac{\partial G(f_{i})}{\partial\kappa_{2}}
&= \frac{1}{\alpha_{2}^{1/2}}\,\frac{\partial G(f_{i})}{\partial P_{2}} \\[2pt]
\frac{\partial G(f_{i})}{\partial L_{2}}
&= \frac{1}{\alpha_{2}^{1/2}}\,\frac{\partial G(f_{i})}{\partial Q_{2}}
\end{aligned}
\qquad ,\ \forall f_{i} .
\]

L'élimination de $\partial G/\partial Q_{2}$ et $\partial G/\partial P_{2}$ fournit la relation centrale
de la note :
\[ 2\alpha_{2}\,\frac{\partial G(f_{i})}{\partial\alpha_{2}}
+ \kappa_{2}\,\frac{\partial G(f_{i})}{\partial\kappa_{2}}
+ L_{2}\,\frac{\partial G(f_{i})}{\partial L_{2}} = 0 ,
\qquad \forall f_{i} . \]

Elle est valable pour toute fréquence et pour toute valeur des trois paramètres, malgré leur influence
non linéaire sur la température. Les trois vecteurs de sensibilité sont donc linéairement dépendants :
le triplet $\{L_{2},\alpha_{2},\kappa_{2}\}$ est pleinement corrélé et non identifiable.

\subsection{Lecture par groupe d'invariance}

La relation précédente est l'identité d'Euler associée au groupe à un paramètre
\[ L_{2}\mapsto\mu\,L_{2} , \qquad \alpha_{2}\mapsto\mu^{2}\alpha_{2} ,
\qquad \kappa_{2}\mapsto\mu\,\kappa_{2} , \qquad \mu>0 . \]

Sous cette action, les grandeurs dérivées se transforment selon
\[ \rho_{2}C_{2}\mapsto\mu^{-1}\rho_{2}C_{2} , \qquad P_{2}\mapsto P_{2} ,
\qquad Q_{2}\mapsto Q_{2} , \]
c'est-à-dire que l'effusivité et le temps de transit demeurent inchangés. Le quadripôle étant fonction
de ces deux seules quantités, la réponse de surface est rigoureusement invariante le long de l'orbite,
pour toute fréquence et en l'absence de bruit.

\begin{center}
\begin{tabular}{@{}lcc@{}}
\toprule
\textbf{Grandeur} & \textbf{Exposant} & \textbf{Statut} \\
\midrule
Épaisseur $L_{2}$ & $\mu^{1}$ & Variable le long de l'orbite \\
Diffusivité $\alpha_{2}$ & $\mu^{2}$ & Variable le long de l'orbite \\
Conductivité $\kappa_{2}$ & $\mu^{1}$ & Variable le long de l'orbite \\
Capacité volumique $\rho_{2}C_{2}$ & $\mu^{-1}$ & Variable le long de l'orbite \\
Effusivité $P_{2}$ & $\mu^{0}$ & Invariante \\
Temps de transit $Q_{2}$ & $\mu^{0}$ & Invariante \\
\bottomrule
\end{tabular}
\end{center}

L'ensemble des solutions admissibles est donc une courbe et non un point. L'incertitude portée par
l'épaisseur n'est pas grande : elle est non bornée.

\subsection{Le jeu $\{L_{2},\alpha_{2},P_{2}\}$}

Le remplacement de la conductivité par l'effusivité ne résout rien. Les sensibilités à la diffusivité et
à l'épaisseur restent linéairement dépendantes :
\[ 2\alpha_{2}\,\frac{\partial G(f_{i})}{\partial\alpha_{2}}
+ L_{2}\,\frac{\partial G(f_{i})}{\partial L_{2}} = 0 , \qquad\forall f_{i} , \]
ce qui annule à nouveau les déterminants correspondants. Le triplet $\{L_{2},\alpha_{2},P_{2}\}$ n'est
pas davantage identifiable.

\subsection{Conséquence sur les cartes d'identifiabilité}

Les cartes de l'article commenté sont censées signaler par des marqueurs les configurations où la
condition d'identifiabilité n'est pas remplie, c'est-à-dire où le déterminant de la matrice des
coefficients de sensibilité s'annule. La note relève deux objections.

\begin{itemize}
\item Le déterminant considéré fait intervenir trois valeurs de fréquence, alors que chaque marqueur des
cartes est rattaché à une fréquence unique.
\item La relation de dépendance linéaire étant valable pour toute fréquence et pour toute valeur des
paramètres, ces déterminants sont identiquement nuls. Une carte correcte serait uniformément couverte de
marqueurs : la condition d'identifiabilité n'est jamais satisfaite.
\end{itemize}

Les valeurs non nulles obtenues par les auteurs sont attribuées aux erreurs d'arrondi accumulées dans
l'approximation par différences finies des dérivées.

\subsection{Conséquence sur la validation expérimentale}

La validation expérimentale présentée dans l'article commenté repose sur l'identification simultanée de
l'épaisseur totale et des propriétés thermophysiques, la première couche étant thermiquement mince et
donc fusionnée avec le revêtement. La note tient cette estimation pour théoriquement impossible, quelle
que soit la méthodologie employée, puisque la matrice d'information est singulière.

Il en résulte que toute autre valeur d'épaisseur située sur l'orbite d'invariance, associée à la
diffusivité correspondante, ajuste les données aussi bien. Les éléments diagonaux de la matrice de
covariance n'ont pas été fournis ; la note observe qu'ils seraient très grands, et que l'affirmation
d'unicité de la mesure est infondée.

\section{Critique des critères d'identifiabilité employés}

\begin{center}
\begin{tabularx}{\textwidth}{@{}lXX@{}}
\toprule
\textbf{Critère} & \textbf{Usage dans l'article commenté} & \textbf{Objection de la note} \\
\midrule
Annulation d'une courbe de sensibilité & Retenue comme signe de non-iden\-ti\-fia\-bi\-li\-té &
L'identification n'utilise pas une fréquence unique mais l'ensemble de la bande \\
\addlinespace
Déterminant de la matrice de sensibilité & Fondé sur deux ou trois fréquences &
Restriction arbitraire alors que les campagnes expérimentales en comportent 19 ou 59 \\
\addlinespace
Déterminant de $S^{T}S$ & Non employé & Critère recommandé ; sa maximisation minimise la région de
confiance \\
\addlinespace
Nombre de conditionnement & Non employé & Signale les jeux de paramètres pathologiques \\
\addlinespace
Diagonale de $\left(S^{T}S\right)^{-1}$ & Non fournie & Donne directement les facteurs d'amplification
de la variance \\
\addlinespace
Sensibilités réduites $\beta\,\partial G/\partial\beta$ & Non employées & Rendent les sensibilités
homogènes et donc comparables entre elles \\
\bottomrule
\end{tabularx}
\end{center}

\section{Ce qui demeure identifiable}

La conclusion opérationnelle de la note est la suivante. Dans la configuration examinée, deux quantités
au plus sont accessibles :

\begin{itemize}
\item le rapport d'effusivité $b_{32}$ ;
\item la racine du temps de transit $Q_{2}$, sous réserve que $b_{32}$ ne soit pas trop proche de
l'unité, la bande de fréquences devant être optimisée pour minimiser les incertitudes.
\end{itemize}

L'épaisseur et la diffusivité ne s'obtiennent ensuite l'une de l'autre qu'au prix d'une mesure
indépendante. Si l'effusivité du substrat est connue, $b_{32}$ fournit celle du revêtement, donc $P_{2}$ ;
l'épaisseur se déduit alors de $Q_{2}$ moyennant la connaissance indépendante de la conductivité ou de
la capacité thermique volumique. L'affirmation selon laquelle l'épaisseur serait identifiable
séparément des propriétés thermophysiques est déclarée fausse.

\section{Portée et limites}

\subsection{Généralité du résultat}

Le raisonnement ne dépend d'aucune particularité du revêtement. Il vaut pour toute couche d'épaisseur
finie au sein d'un empilement multicouche, puisque la contribution de chaque couche passe par un
quadripôle de même structure.

\subsection{Voies de levée de l'indétermination}

La note mentionne deux dispositifs qui apportent une information de nature différente et brisent
l'invariance :

\begin{itemize}
\item une expérience photothermique bidimensionnelle mettant en jeu la diffusion latérale, et non
seulement la diffusion transverse ;
\item une expérience comportant des sources internes au sein de la couche étudiée.
\end{itemize}

Dans les deux cas, une longueur physique distincte de l'épaisseur entre dans le modèle direct, ce qui
détruit l'invariance d'échelle décrite en 5.3. Ces configurations sortent du cadre de l'article
commenté.

\subsection{Ce que la note n'aborde pas}

Le cadre constitutif reste celui de la loi de Fourier. Aucune analyse spectrale de l'opérateur associé
n'est entreprise, et la question du conditionnement, distincte de celle de l'identifiabilité, n'est
traitée que par renvoi à des critères existants.

\section{Rattachement au présent travail}

\subsection{Continuité formelle}

Le couple $\{Q_{2},P_{2}\}$ est le couple $\{\xi_{1},b\}$ des conventions du projet : le quadripôle de
l'équation (1) de la note est celui du mur homogène exprimé dans la coordonnée de Liouville. La
vérification terme à terme déjà consignée dans le fichier de conventions s'applique sans modification.

\subsection{Statut de l'argument}

La note établit un résultat de non-unicité structurelle, indépendant du bruit et de la bande de mesure.
Elle relève donc de la même catégorie que l'indétermination de la transformation inverse de Liouville,
et se distingue du mauvais conditionnement de l'inversion en profondeur. La distinction entre les deux
sources de non-détermination, déjà posée dans le présent travail, trouve ici un cas d'espèce
élémentaire, où l'orbite d'invariance est explicite et de dimension un.

\subsection{Point de vigilance}

L'invariance décrite en 5.3 concerne le cas homogène. Le cas gradué comporte une invariance de même
nature mais de dimension fonctionnelle, obtenue par reparamétrage de la coordonnée de profondeur. Ce
prolongement fait l'objet de la fiche relative à la controverse de 2023.

\section{Références}

\begin{enumerate}[leftmargin=2.2em]
\item J.-C. Krapez, F. Rigollet, \emph{Comment on ``Photothermal radiometry parametric identifiability
theory for reliable and unique nondestructive coating thickness and thermophysical measurements''
[J. Appl. Phys. 121(9), 095101 (2017)]}, arXiv:1708.07362v1, 2017.
\item X. Guo, A. Mandelis, J. Tolev, K. Tang, \emph{Journal of Applied Physics} \textbf{121}, 095101
(2017).
\item D. Maillet, S. André, J.-C. Batsale, A. Degiovanni, C. Moyne, \emph{Thermal Quadrupoles: Solving
the Heat Equation through Integral Transforms}, Wiley, 2000.
\item J. Beck, K. Arnold, \emph{Parameter Estimation in Engineering and Science}, Wiley, 1977.
\item H. S. Carslaw, J. C. Jaeger, \emph{Conduction of Heat in Solids}, Oxford University Press, 1959.
\end{enumerate}
\end{document}
